\exercisesection{3}

¶ 1. (a) 设 L 为域，T 为不定元；在形式幂级数域 L((T)) 中，考察级数

\[
s = c _ {0} + c _ {1} \mathrm{T} ^ {1!} + c _ {2} \mathrm{T} ^ {2!} + \dots + c _ {n} \mathrm{T} ^ {n!} + \dots
\]

其中 \(c_{i} \in L\) 均 \(\neq 0\)。证明 s 在 \(\mathrm{L}(\mathrm{T})\) 上超越；\(\mathrm{L}((\mathrm{T}))\) 是其扩域。（反证，假设存在关系 \(a_{0}(\mathrm{T}) + a_{1}(\mathrm{T})s + \cdots + a_{q}(\mathrm{T})s^{q} = 0\)，其中 \(a_{i}(\mathrm{T})\) 是 \(\mathrm{L}(\mathrm{T})\) 中的多项式且 \(a_{q} \neq 0\)；构造下述元素所满足的方程

\[
s ^ {\prime} = s - c _ {0} - c _ {1} \mathrm{T} ^ {1!} - \dots - c _ {n - 1} \mathrm{T} ^ {(n - 1)!}
\]

并证明当 \(n\) 足够大时，其常数项是次数 \(< n!\) 的多项式，从而矛盾。）

\begin{enumerate}
\def\labelenumi{(\alph{enumi})}
\setcounter{enumi}{1}
\tightlist
\item
  设 \(k\) 为域，\(\mathbf{X}, \mathbf{Y}, \mathbf{Z}\)、\(\mathrm{T}\) 为四个不定元，\(\mathbf{K} = k(\mathbf{X}, \mathbf{Y}, \mathbf{Z})\)，\(\mathbf{E}\) 为 \(k(\mathbf{X})\) 的代数闭包；在 \(\mathbf{E}\) 中，对所有 \(n\)，令 \(x_{n}\) 表示满足 \(x_{n}^{n} = \mathbf{X}\) 的元素。在形式幂级数域 \(\mathrm{E}((\mathrm{T}))\) 中，考察元素
\end{enumerate}

\[
s = x _ {1} \mathrm{T} + x _ {2} \mathrm{T} ^ {2!} + \dots + x _ {n} \mathrm{T} ^ {n!} + \dots
\]

由 (a)，元素 \(s\) 和 \(\mathbf{T}\) 在 \(\mathbf{E}\) 上代数独立。考察同态 \(f\colon \mathbf{K} \to \mathbf{E}((\mathbf{T}))\)，满足 \(f(\mathbf{X}) = \mathbf{X}, f(\mathbf{Y}) = \mathbf{T}, f(\mathbf{Z}) = s\)。若 \(v\) 是 \(\mathbf{E}((\mathbf{T}))\) 上由形式幂级数的阶定义的离散赋值（第 3 节例 3），则 \(v \circ f\) 是 \(\mathbf{K}\) 上的离散赋值；证明 \(v \circ f\) 的剩余域是子域 \(k(x_1, x_2, \ldots, x_n, \ldots)\)，它属于 \(\mathbf{E}\)，因而在 \(k\) 上是无限的。

\begin{enumerate}
\def\labelenumi{\arabic{enumi}.}
\setcounter{enumi}{1}
\tightlist
\item
  设 \(\Gamma\) 为全序群，\(k\) 为域。
\end{enumerate}

\begin{enumerate}
\def\labelenumi{(\alph{enumi})}
\item
  设 A、B 是 \(\Gamma\) 的两个子集，按诱导序良序；证明集合 \(A + B\) 是良序的，并且对于每个 \(\gamma \in A + B\)，满足 \((\alpha, \beta) \in A \times B\) 且 \(\alpha + \beta = \gamma\) 的有序对的集合是有限的。（为证明第一个断言，反证：否则可以定义严格递减序列 \((\gamma_n)\)，其各项属于 \(A + B\)，使 \(\gamma_n = \alpha_n + \beta_n\)，其中 \(\alpha_n \in A\)、\(\beta_n \in B\)，序列 \((\alpha_n)\) 严格递增而序列 \((\beta_n)\) 严格递减；使用《集合论》第三章 § 4 习题 3 得出结论。）
\item
  设 \(S(\Gamma, k)\) 为下述族的集合：\(x = (x_{\alpha})_{\alpha \in \Gamma}\)，满足 \(x_{\alpha} \in k\) 对所有 \(\alpha \in \Gamma\) 成立，且集合 \(\alpha\)（满足 \(x_{\alpha} \neq 0\)）是 \(\Gamma\) 的良序子集。证明 \(S(\Gamma, k)\) 是 \(k^{\Gamma}\) 的加法子群；此外，若 \(x = (x_{\alpha}), y = (y_{\alpha})\) 是 \(S(\Gamma, k)\) 中的两个元素，则集合 \(C\) 由 \(\alpha + \beta\) 组成，其中有序对 \((\alpha, \beta)\) 满足 \(x_{\alpha} \neq 0\) 且 \(y_{\beta} \neq 0\)；由 (a)，它是良序的，并且对每个 \(\gamma \in C\)，元素 \(z_{\gamma} = \sum_{\alpha + \beta = \gamma} x_{\alpha} y_{\beta}\) 在 \(k\) 中有定义；若置 \(z_{\alpha} = 0\) 对 \(\alpha \notin C\)，则元素 \(z = (z_{\alpha}) \in k^{\Gamma}\) 属于 \(S(\Gamma, k)\)；将其记为 \(xy\)。证明在乘积群 \(k^{\Gamma}\) 上赋予上述乘法和加法后，集合 \(S(\Gamma, k)\) 是域；此外，对于 \(x = (x_{\alpha}) \neq 0\) 属于 \(S(\Gamma, k)\) 的任意元素，令 \(\lambda\) 为 \(\alpha \in \Gamma\) 中使 \(x_{\alpha} \neq 0\) 的最小者；若记 \(v(x) = \lambda\)（并令 \(v(0) = +\infty\)），证明 \(v\) 是 \(S(\Gamma, k)\) 上的赋值，剩余域为 \(k\)，值群为 \(\Gamma\)。称 \(v\) 为域 \(S(\Gamma, k)\) 上的典范赋值。域 \(k\) 典范地识别为 \(S(\Gamma, k)\) 的一个子域。对所有 \(\alpha \in \Gamma\)，令 \(u_{\alpha}\) 表示元素 \((x_{\lambda})_{\lambda \in \Gamma}\)，满足 \(x_{\alpha} = 1, x_{\lambda} = 0\) 对 \(\lambda \neq \alpha\)；于是，对每个非零 \(z \in S(\Gamma, k)\)，存在唯一的有序对 \((t, \alpha) \in k \times \Gamma\)，使 \(v(z) = \alpha\) 且 \(v(z - tu_{\alpha}) > \alpha\)。称 \(tu_{\alpha}\) 为 \(z\) 的主导项。
\end{enumerate}

¶ 3. 设 K 为域，w 为 K 上的赋值，Γ 为 w 的值群，k 为 w 的剩余域。构造相应的域 S(Γ, k)，并沿用习题 2 的记号。对于所有 \(t \in k\)，令 \(t^{0}\) 为 w 的赋值环中类 t 的一个元素；对于所有 \(\alpha \in \Gamma\)，令 \(u_{\alpha}^{0}\) 为 K 中满足 \(w(u_{\alpha}^{0}) = \alpha\) 的元素。

\begin{enumerate}
\def\labelenumi{(\alph{enumi})}
\tightlist
\item
  设 M 是 S(Γ, k) 的一个子集，包含元素 \(tu_{\alpha}\)，其中每个有序对 \((t, \alpha) \in k \times \Gamma\) 都对应一个这样的元素；假设存在双射 \(x \mapsto x^{0}\)，它把 M 映到 K 的一个子集 \(M^{0}\)，且满足下列条件：(1) \((tu_{\alpha})^{0} = t^{0}u_{\alpha}^{0}\) 对每个有序对 \((t, \alpha) \in k \times \Gamma\) 成立；(2) 若 \(x = (x_{\alpha})\) 属于 M，且 \(C_{x}\) 是 Γ 的良序子集，由满足 \(x_{\alpha} \neq 0\) 的 α 组成，则对 \(C_{x}\) 的每个区段，元素 \((x_{\lambda})_{\lambda \in D}\) 属于 M；(3) 若 x、y 是 M 的两个元素，且 \(tu_{\alpha}\) 是 x - y 的主导项，则 \(w(x^{0} - y^{0} - t^{0}u_{\alpha}^{0}) > w(x^{0} - y^{0})\)。若 \(z'\) 是 K 中不属于 \(M^{0}\) 的元素，证明存在 S(Γ, k) 中的元素 \(z \notin M\)，使得在 M 上与 \(x \mapsto x^{0}\) 重合且把 \(z'\) 映到 z 的映射也满足上述条件。（注意，若 \(w(z') = \alpha\)，存在 \(t^{0}\)，使
\end{enumerate}

\[
w (z ^ {\prime} - t ^ {0} u _ {\alpha} ^ {0}) > \alpha ;
\]

令 \(z_{\alpha} = tu_{\alpha}\)，然后证明通过超限归纳可以确定 \(z_{\beta}\)（其中 \(\beta > \alpha\)），使 \(z = (z_{\lambda})_{\lambda \in \Gamma}\) 解决此问题。）

\begin{enumerate}
\def\labelenumi{(\alph{enumi})}
\setcounter{enumi}{1}
\tightlist
\item
  由 (a) 推出 \(\operatorname{Card}(\mathbf{K}) \leqslant \operatorname{Card} S(\Gamma, k)\)。
\end{enumerate}

\begin{enumerate}
\def\labelenumi{\arabic{enumi}.}
\setcounter{enumi}{3}
\tightlist
\item
  设 A 是局部整环，K 是其分式域，m 是其极大理想，v 是 K 上的赋值，且 v 的环 B 支配 A。假设 m 是有限生成理想，或者 v 是离散赋值；则 m 中存在元素 x，使 \(v(x) = \inf_{y \in \mathfrak{m}} v(y)\)。令 \(A'\) 为 K 中由
\end{enumerate}

A 以及 \(yx^{-1}\)（其中 y 遍历 m）生成的子环；令 \(A_{1}\) 为 \(A'\) 的局部环，相对于素理想 \(p \cap A'\)，其中 p 是赋值 v 的理想。证明环 \(A_{1}\) 不依赖于满足上述条件的元素 x 的选取（\(A_{1}\) 称为“在 \(v'\) 方向上 A 的素二次变换”）；若 A 是 Noether 环，则 \(A_{1}\) 也是 Noether 环。

¶ 5. 设 k 为域，K = k(X, Y) 是 k 上两个不定元的有理函数域。记 \(x_{1} = X\)、\(y_{1} = Y\)，并对 \(n \geqslant 1\) 归纳定义 K 中的元素 \(x_{n+1} = y_{n}\) 和 \(y_{n+1} = x_{n}y_{n}^{-1}\)。

\begin{enumerate}
\def\labelenumi{(\alph{enumi})}
\tightlist
\item
  写 \(\mathbf{A}_n = k[x_n, y_n]\)；序列 \((\mathbf{A}_n)\) 递增，并且若
\end{enumerate}

\[
\mathfrak {m} _ {n} = \mathrm{A} _ {n} x _ {n} + \mathrm{A} _ {n} y _ {n},
\]

\(m_{n}\) 是 \(A_{n}\) 的极大理想，且 \(m_{n} = A_{n} \cap m_{n+1}\)；在环 \(A = \bigcap_{n} A_{n}\) 中，\(m = \bigcap_{n} m_{n}\) 是极大理想。

\begin{enumerate}
\def\labelenumi{(\alph{enumi})}
\setcounter{enumi}{1}
\item
  设 \(v\) 为 \(\mathbf{K}\) 上由 \(v(\mathbf{X}) = \rho = (1 + \sqrt{5}) / 2\)、\(v(\mathbf{Y}) = 1\) 定义的赋值；令 \(\mathbf{B}\) 和 \(\mathfrak{p}\) 为 \(v\) 的环和理想；证明 \(\mathfrak{m} = \mathfrak{p} \cap \mathbf{A}\)（注意 \(\rho^2 - \rho - 1 = 0\)，并用此计算 \(v(y_n)\)）。证明对于单项式 \(\mathbf{X}^i\mathbf{Y}^j\)（\(i, j\) 为正整数或负整数）属于 \(\mathbf{A}\) 的充要条件是 \(i\left(\frac{1 + \sqrt{5}}{2}\right) + j > 0\)；因此，对于每个单项式 \(z = \mathbf{X}^i\mathbf{Y}^j\)，要么 \(z \in \mathbf{A}\)，要么 \(1/z \in \mathbf{A}\)。
\item
  证明 \(B = A_{m}\)（将 \(t \in B\) 写成 X、Y 中两个多项式之商的形式，并证明在这两个多项式中，v 取最小值的单项式……）。
\item
  证明对所有 \(n\)，\((\mathbf{A}_{n+1})_{\mathfrak{m}_{n+1}}\) 是 \((\mathbf{A}_n)_{\mathfrak{m}_n}\) 在 \(v\) 方向上的素二次变换（练习 4）。
\end{enumerate}

\begin{enumerate}
\def\labelenumi{\arabic{enumi}.}
\setcounter{enumi}{5}
\tightlist
\item
  \begin{enumerate}
  \def\labelenumii{(\alph{enumii})}
  \tightlist
  \item
    设 \(\mathbf{A}\) 为赋值环。将《代数》第七章 § 4 第 1 至 4 节的结果（不变因子理论）推广到有限生成 A-模。
  \end{enumerate}
\end{enumerate}

\begin{enumerate}
\def\labelenumi{(\alph{enumi})}
\setcounter{enumi}{1}
\item
  设 K 为域，v 为 K 上的赋值，A 为 v 的环，且 \(\Gamma\) 为 v 的阶群。令 \(U = (\alpha_{ij})\) 为 A 中元素构成的 n 阶方阵，且其行列式为 0；证明存在 \(\lambda \in \Gamma\)，使得对于每个由 A 中元素构成的 n 阶方阵 \(U' = (\alpha_{ij}')\)，若 \(v(\alpha_{ij}' - \alpha_{ij}) > \lambda\) 对每个有序对 \((i,j)\) 都成立，则 \(U'\) 的不变因子与 U 的相同。
\item
  在 (b) 的假设下，推广《代数》第九章 § 5 第 1 节的定理和习题 1 以及 § 6 习题 10 的下列结果（对于最后一个结果，假设 K 的特征不为 2 且 \(v(2) = 0\)）。
\end{enumerate}

\begin{enumerate}
\def\labelenumi{\arabic{enumi}.}
\setcounter{enumi}{6}
\item
  证明第二章 § 3 习题 2 (b) 中定义的整环 B 是局部环，其极大理想是主理想，但 B 不是赋值环。
\item
  证明若赋值环 A 是强 Lasker 环（第四章 § 2 习题 28），则 A 是域或离散赋值环（使用第四章 § 2 习题 29）。
\end{enumerate}
% semantic-fragments: legacy-input-seam
\relax{}
